
\documentclass{article}
\usepackage{colortbl}
\usepackage{makecell}
\usepackage{multirow}
\usepackage{supertabular}

\begin{document}

\newcounter{utterance}

\centering \large Interaction Transcript for game `hot\_air\_balloon', experiment `air\_balloon\_survival\_en\_reasoning off\_hard', episode 1 with gpt{-}5{-}2025{-}08{-}07{-}t1.0.
\vspace{24pt}

{ \footnotesize  \setcounter{utterance}{1}
\setlength{\tabcolsep}{0pt}
\begin{supertabular}{c@{$\;$}|p{.15\linewidth}@{}p{.15\linewidth}p{.15\linewidth}p{.15\linewidth}p{.15\linewidth}p{.15\linewidth}}
    \# & $\;$A & \multicolumn{4}{c}{Game Master} & $\;\:$B\\
    \hline

    \theutterance \stepcounter{utterance}  
    & & \multicolumn{4}{p{0.6\linewidth}}{
        \cellcolor[rgb]{0.9,0.9,0.9}{
            \makecell[{{p{\linewidth}}}]{
                \texttt{\tiny{[P1$\langle$GM]}}
                \texttt{Stai partecipando a un gioco di negoziazione collaborativa.} \\
\\ 
\texttt{Insieme a un altro partecipante, dovete accordarvi su un unico set di oggetti da conservare. Ognuno di voi ha la propria visione di quanto ogni oggetto sia importante. Non sai come l'altro partecipante valuta gli oggetti. Inoltre, ti viene fornito lo sforzo che ogni oggetto richiede.} \\
\texttt{Potete accordarvi su un set solo se lo sforzo totale degli oggetti selezionati non supera un limite condiviso:} \\
\\ 
\texttt{LIMITE = 3357} \\
\\ 
\texttt{Ecco i valori di sforzo dei singoli oggetti:} \\
\\ 
\texttt{Sforzo oggetto = \{"A21": 506, "B93": 492, "A11": 338, "B33": 122, "C08": 131, "B39": 909, "B96": 144, "C85": 716, "A82": 263, "B31": 231, "B20": 91, "A60": 651, "B09": 552, "C81": 849, "B21": 720\}} \\
\\ 
\texttt{Ecco la tua visione personale sull'importanza di ogni oggetto:} \\
\\ 
\texttt{Valori di importanza dell'oggetto = \{"A21": 138, "B93": 583, "A11": 868, "B33": 822, "C08": 783, "B39": 65, "B96": 262, "C85": 121, "A82": 508, "B31": 780, "B20": 461, "A60": 484, "B09": 668, "C81": 389, "B21": 808\}} \\
\\ 
\texttt{Obiettivo:} \\
\\ 
\texttt{Il tuo obiettivo è negoziare un set condiviso di oggetti che ti avvantaggi il più possibile (cioè, massimizzi l'importanza totale per TE), rimanendo entro il LIMITEE. Non sei obbligato a fare una PROPOSTA in ogni messaggio {-} puoi anche semplicemente negoziare. Tutte le tattiche sono permesse!} \\
\\ 
\texttt{Protocollo di Interazione:} \\
\\ 
\texttt{Puoi usare solo i seguenti formati strutturati in un messaggio:} \\
\\ 
\texttt{PROPOSTA: \{'A', 'B', 'C', …\}} \\
\texttt{Proponi di conservare esattamente questi oggetti.} \\
\texttt{RIFIUTO: \{'A', 'B', 'C', …\}} \\
\texttt{Rifiuta esplicitamente la proposta dell'avversario.} \\
\texttt{ARGOMENTAZIONE: \{'...'\}} \\
\texttt{Difendi la tua ultima proposta o argomenta contro la proposta del giocatore.} \\
\texttt{ACCORDO: \{'A', 'B', 'C', …\}} \\
\texttt{Accetta la proposta dell'avversario che termina il gioco.} \\
\\ 
\\ 
\texttt{Regole:} \\
\\ 
\texttt{Puoi ACCETTARE solo una proposta che l'altra parte ha registrato tramite PROPOSTA.} \\
\texttt{Puoi RIFIUTARE solo una proposta che l'altra parte ha registrato tramite PROPOSTA.} \\
\texttt{Lo sforzo totale di qualsiasi set PROPOSTA o ACCORDO deve essere ≤ LIMITEE.} \\
\texttt{NON rivelare i tuoi punteggi di importanza nascosti.} \\
\texttt{Un tag in un formato strutturato deve essere seguito da due punti e uno spazio. L'argomento deve essere un set python contenente 0 o più stringhe.} \\
\texttt{Quindi, deve essere della forma TAG: \{...\}} \\
\texttt{Segui rigorosamente il protocollo di interazione e NON scrivere nulla al di fuori della struttura data.} \\
\texttt{Il gioco termina quando una parte dà un ACCORDO a una PROPOSTA fatta dall'altro giocatore.} \\
\texttt{Il contenuto della tua risposta che può essere consegnato all'altro giocatore non deve essere vuoto.} \\
\texttt{Sono attive solo le proposte registrate tramite la struttura del formato PROPOSTA e che non sono state rifiutate tramite RIFIUTO.} \\
\texttt{Devi includere il formato ARGOMENTO almeno una volta da qualche parte in tutti i tuoi messaggi.} \\
            }
        }
    }
    & \\ \\

    \theutterance \stepcounter{utterance}  
    & \multicolumn{4}{p{0.6\linewidth}}{
        \cellcolor[rgb]{0.9,0.9,0.9}{
            \makecell[{{p{\linewidth}}}]{
                \texttt{\tiny{[P1$\rangle$GM]}}
                \texttt{ARGOMENTAZIONE: \{'Propongo una selezione ad alta efficienza di sforzo: molti elementi leggeri con impatto trasversale, evitando quelli molto pesanti o poco incisivi. Rimangono 273 di margine per eventuali scambi. Se desiderate inserire un escluso, proponete una sostituzione di sforzo simile mantenendo il limite.'\}} \\
\texttt{PROPOSTA: \{'B33', 'C08', 'B20', 'B31', 'A11', 'A82', 'B96', 'B93', 'B09', 'B21'\}} \\
            }
        }
    }
    & & \\ \\

    \theutterance \stepcounter{utterance}  
    & & \multicolumn{4}{p{0.6\linewidth}}{
        \cellcolor[rgb]{0.9,0.9,0.9}{
            \makecell[{{p{\linewidth}}}]{
                \texttt{\tiny{[GM$\rangle$P2]}}
                \texttt{Stai partecipando a un gioco di negoziazione collaborativa.} \\
\\ 
\texttt{Insieme a un altro partecipante, dovete accordarvi su un unico set di oggetti da conservare. Ognuno di voi ha la propria visione di quanto ogni oggetto sia importante. Non sai come l'altro partecipante valuta gli oggetti. Inoltre, ti viene fornito lo sforzo che ogni oggetto richiede.} \\
\texttt{Potete accordarvi su un set solo se lo sforzo totale degli oggetti selezionati non supera un limite condiviso:} \\
\\ 
\texttt{LIMITE = 3357} \\
\\ 
\texttt{Ecco i valori di sforzo dei singoli oggetti:} \\
\\ 
\texttt{Sforzo oggetto = \{"A21": 506, "B93": 492, "A11": 338, "B33": 122, "C08": 131, "B39": 909, "B96": 144, "C85": 716, "A82": 263, "B31": 231, "B20": 91, "A60": 651, "B09": 552, "C81": 849, "B21": 720\}} \\
\\ 
\texttt{Ecco la tua visione personale sull'importanza di ogni oggetto:} \\
\\ 
\texttt{Valori di importanza dell'oggetto = \{"A11": 65, "B33": 121, "B21": 138, "C08": 262, "B31": 389, "B09": 461, "B93": 484, "A82": 508, "A60": 583, "B20": 668, "C81": 780, "B96": 783, "A21": 808, "C85": 822, "B39": 868\}} \\
\\ 
\texttt{Obiettivo:} \\
\\ 
\texttt{Il tuo obiettivo è negoziare un set condiviso di oggetti che ti avvantaggi il più possibile (cioè, massimizzi l'importanza totale per TE), rimanendo entro il LIMITEE. Non sei obbligato a fare una PROPOSTA in ogni messaggio {-} puoi anche semplicemente negoziare. Tutte le tattiche sono permesse!} \\
\\ 
\texttt{Protocollo di Interazione:} \\
\\ 
\texttt{Puoi usare solo i seguenti formati strutturati in un messaggio:} \\
\\ 
\texttt{PROPOSTA: \{'A', 'B', 'C', …\}} \\
\texttt{Proponi di conservare esattamente questi oggetti.} \\
\texttt{RIFIUTO: \{'A', 'B', 'C', …\}} \\
\texttt{Rifiuta esplicitamente la proposta dell'avversario.} \\
\texttt{ARGOMENTAZIONE: \{'...'\}} \\
\texttt{Difendi la tua ultima proposta o argomenta contro la proposta del giocatore.} \\
\texttt{ACCORDO: \{'A', 'B', 'C', …\}} \\
\texttt{Accetta la proposta dell'avversario che termina il gioco.} \\
\\ 
\\ 
\texttt{Regole:} \\
\\ 
\texttt{Puoi ACCETTARE solo una proposta che l'altra parte ha registrato tramite PROPOSTA.} \\
\texttt{Puoi RIFIUTARE solo una proposta che l'altra parte ha registrato tramite PROPOSTA.} \\
\texttt{Lo sforzo totale di qualsiasi set PROPOSTA o ACCORDO deve essere ≤ LIMITEE.} \\
\texttt{NON rivelare i tuoi punteggi di importanza nascosti.} \\
\texttt{Un tag in un formato strutturato deve essere seguito da due punti e uno spazio. L'argomento deve essere un set python contenente 0 o più stringhe.} \\
\texttt{Quindi, deve essere della forma TAG: \{...\}} \\
\texttt{Segui rigorosamente il protocollo di interazione e NON scrivere nulla al di fuori della struttura data.} \\
\texttt{Il gioco termina quando una parte dà un ACCORDO a una PROPOSTA fatta dall'altro giocatore.} \\
\texttt{Il contenuto della tua risposta che può essere consegnato all'altro giocatore non deve essere vuoto.} \\
\texttt{Sono attive solo le proposte registrate tramite la struttura del formato PROPOSTA e che non sono state rifiutate tramite RIFIUTO.} \\
\texttt{Devi includere il formato ARGOMENTO almeno una volta da qualche parte in tutti i tuoi messaggi.} \\
\\ 
\\ 
\texttt{Ora riceverai il primo messaggio dell'altro giocatore.} \\
\\ 
\texttt{ARGOMENTAZIONE: \{'Propongo una selezione ad alta efficienza di sforzo: molti elementi leggeri con impatto trasversale, evitando quelli molto pesanti o poco incisivi. Rimangono 273 di margine per eventuali scambi. Se desiderate inserire un escluso, proponete una sostituzione di sforzo simile mantenendo il limite.'\}} \\
\texttt{PROPOSTA: \{'B33', 'C08', 'B20', 'B31', 'A11', 'A82', 'B96', 'B93', 'B09', 'B21'\}} \\
            }
        }
    }
    & \\ \\

    \theutterance \stepcounter{utterance}  
    & & & \multicolumn{4}{p{0.6\linewidth}}{
        \cellcolor[rgb]{0.9,0.9,0.9}{
            \makecell[{{p{\linewidth}}}]{
                \texttt{\tiny{[GM$\langle$P2]}}
                \texttt{ARGOMENTAZIONE: \{'Accolgo il tuo criterio di efficienza e mantengo 7 su 10 delle tue scelte. Propongo tre sostituzioni mirate usando il buffer senza superare il limite: A11 {-}> A21 (+168), B21 {-}> A60 ({-}69), B09 {-}> C85 (+164). Sforzo totale risultante: 3347, margine 10. In questo modo rimuoviamo un pezzo pesante e poco trasversale e inseriamo due elementi ad alto impatto per unita di sforzo, preservando agilita complessiva. Se una di queste sostituzioni e critica per te, indicala e posso proporre una variante equivalente a sforzo simile.'\}} \\
\texttt{PROPOSTA: \{'B33', 'C08', 'B20', 'B31', 'A21', 'A82', 'B96', 'B93', 'C85', 'A60'\}} \\
            }
        }
    }
     \\ \\

    \theutterance \stepcounter{utterance}  
    & & \multicolumn{4}{p{0.6\linewidth}}{
        \cellcolor[rgb]{0.9,0.9,0.9}{
            \makecell[{{p{\linewidth}}}]{
                \texttt{\tiny{[P1$\langle$GM]}}
                \texttt{ARGOMENTAZIONE: \{'Accolgo il tuo criterio di efficienza e mantengo 7 su 10 delle tue scelte. Propongo tre sostituzioni mirate usando il buffer senza superare il limite: A11 {-}> A21 (+168), B21 {-}> A60 ({-}69), B09 {-}> C85 (+164). Sforzo totale risultante: 3347, margine 10. In questo modo rimuoviamo un pezzo pesante e poco trasversale e inseriamo due elementi ad alto impatto per unita di sforzo, preservando agilita complessiva. Se una di queste sostituzioni e critica per te, indicala e posso proporre una variante equivalente a sforzo simile.'\}} \\
\texttt{PROPOSTA: \{'B33', 'C08', 'B20', 'B31', 'A21', 'A82', 'B96', 'B93', 'C85', 'A60'\}} \\
            }
        }
    }
    & \\ \\

    \theutterance \stepcounter{utterance}  
    & \multicolumn{4}{p{0.6\linewidth}}{
        \cellcolor[rgb]{0.9,0.9,0.9}{
            \makecell[{{p{\linewidth}}}]{
                \texttt{\tiny{[P1$\rangle$GM]}}
                \texttt{RIFIUTO: \{'B33', 'C08', 'B20', 'B31', 'A21', 'A82', 'B96', 'B93', 'C85', 'A60'\}} \\
\texttt{ARGOMENTAZIONE: \{'Accolgo l’idea di efficienza. Mantengo A60 come da tua preferenza, ma reintegro A11 e B09 perché offrono copertura ampia con minor sforzo ed evitano sovrapposizioni rispetto a C85 e A21, che aggiungono carico senza benefici proporzionati nel contesto già coperto da C08/B33/B31. Il set proposto pesa 3015 (margine 342), massimizzando agilità. Se una delle due esclusioni è per te critica, posso offrire varianti a sforzo simile: (i) inserire A21 al posto di A11 (totale \textasciitilde{}3183), oppure (ii) inserire C85 al posto di B09 (totale \textasciitilde{}3179). Fammi sapere quale preferisci esplorare.'\}} \\
\texttt{PROPOSTA: \{'B33', 'C08', 'B20', 'B31', 'A11', 'A82', 'B96', 'B93', 'B09', 'A60'\}} \\
            }
        }
    }
    & & \\ \\

    \theutterance \stepcounter{utterance}  
    & & \multicolumn{4}{p{0.6\linewidth}}{
        \cellcolor[rgb]{0.9,0.9,0.9}{
            \makecell[{{p{\linewidth}}}]{
                \texttt{\tiny{[GM$\rangle$P2]}}
                \texttt{RIFIUTO: \{'B33', 'C08', 'B20', 'B31', 'A21', 'A82', 'B96', 'B93', 'C85', 'A60'\}} \\
\texttt{ARGOMENTAZIONE: \{'Accolgo l’idea di efficienza. Mantengo A60 come da tua preferenza, ma reintegro A11 e B09 perché offrono copertura ampia con minor sforzo ed evitano sovrapposizioni rispetto a C85 e A21, che aggiungono carico senza benefici proporzionati nel contesto già coperto da C08/B33/B31. Il set proposto pesa 3015 (margine 342), massimizzando agilità. Se una delle due esclusioni è per te critica, posso offrire varianti a sforzo simile: (i) inserire A21 al posto di A11 (totale \textasciitilde{}3183), oppure (ii) inserire C85 al posto di B09 (totale \textasciitilde{}3179). Fammi sapere quale preferisci esplorare.'\}} \\
\texttt{PROPOSTA: \{'B33', 'C08', 'B20', 'B31', 'A11', 'A82', 'B96', 'B93', 'B09', 'A60'\}} \\
            }
        }
    }
    & \\ \\

    \theutterance \stepcounter{utterance}  
    & & & \multicolumn{4}{p{0.6\linewidth}}{
        \cellcolor[rgb]{0.9,0.9,0.9}{
            \makecell[{{p{\linewidth}}}]{
                \texttt{\tiny{[GM$\langle$P2]}}
                \texttt{ARGOMENTAZIONE: \{'Accolgo la tua impostazione agile. Seguo la variante (i) che hai proposto: sostituiamo A11 con A21. L’incremento di sforzo è di soli +168 (totale 3183, margine 174), restando pienamente entro il limite e senza introdurre sovrapposizioni rispetto a C08/B33/B31. A21 aumenta la copertura e la resilienza del set mantenendo l’efficienza complessiva. Se preferisci invece esplorare l’alternativa (ii), posso valutare C85 al posto di B09 a pari vincoli.'\}} \\
\texttt{PROPOSTA: \{'B33', 'C08', 'B20', 'B31', 'A21', 'A82', 'B96', 'B93', 'B09', 'A60'\}} \\
            }
        }
    }
     \\ \\

    \theutterance \stepcounter{utterance}  
    & & \multicolumn{4}{p{0.6\linewidth}}{
        \cellcolor[rgb]{0.9,0.9,0.9}{
            \makecell[{{p{\linewidth}}}]{
                \texttt{\tiny{[P1$\langle$GM]}}
                \texttt{ARGOMENTAZIONE: \{'Accolgo la tua impostazione agile. Seguo la variante (i) che hai proposto: sostituiamo A11 con A21. L’incremento di sforzo è di soli +168 (totale 3183, margine 174), restando pienamente entro il limite e senza introdurre sovrapposizioni rispetto a C08/B33/B31. A21 aumenta la copertura e la resilienza del set mantenendo l’efficienza complessiva. Se preferisci invece esplorare l’alternativa (ii), posso valutare C85 al posto di B09 a pari vincoli.'\}} \\
\texttt{PROPOSTA: \{'B33', 'C08', 'B20', 'B31', 'A21', 'A82', 'B96', 'B93', 'B09', 'A60'\}} \\
            }
        }
    }
    & \\ \\

    \theutterance \stepcounter{utterance}  
    & \multicolumn{4}{p{0.6\linewidth}}{
        \cellcolor[rgb]{0.9,0.9,0.9}{
            \makecell[{{p{\linewidth}}}]{
                \texttt{\tiny{[P1$\rangle$GM]}}
                \texttt{RIFIUTO: \{'B33', 'C08', 'B20', 'B31', 'A21', 'A82', 'B96', 'B93', 'B09', 'A60'\}} \\
\texttt{ARGOMENTAZIONE: \{'Resto sulla linea “agile” che entrambi abbiamo condiviso. A21 e A60 introducono carico significativo con benefici marginali nel perimetro già coperto da C08/B33/B31. Propongo di sostituirli con A11 e B21: manteniamo copertura ampia senza sovrapposizioni, alta efficienza per unità di sforzo e un margine di 273 (sforzo totale 3084). Se A60 per te è imprescindibile, posso valutare una variante che lo includa mantenendo il limite; segnalamelo e la formalizzo.'\}} \\
\texttt{PROPOSTA: \{'B33', 'C08', 'B20', 'B31', 'A11', 'A82', 'B96', 'B93', 'B09', 'B21'\}} \\
            }
        }
    }
    & & \\ \\

    \theutterance \stepcounter{utterance}  
    & & \multicolumn{4}{p{0.6\linewidth}}{
        \cellcolor[rgb]{0.9,0.9,0.9}{
            \makecell[{{p{\linewidth}}}]{
                \texttt{\tiny{[GM$\rangle$P2]}}
                \texttt{RIFIUTO: \{'B33', 'C08', 'B20', 'B31', 'A21', 'A82', 'B96', 'B93', 'B09', 'A60'\}} \\
\texttt{ARGOMENTAZIONE: \{'Resto sulla linea “agile” che entrambi abbiamo condiviso. A21 e A60 introducono carico significativo con benefici marginali nel perimetro già coperto da C08/B33/B31. Propongo di sostituirli con A11 e B21: manteniamo copertura ampia senza sovrapposizioni, alta efficienza per unità di sforzo e un margine di 273 (sforzo totale 3084). Se A60 per te è imprescindibile, posso valutare una variante che lo includa mantenendo il limite; segnalamelo e la formalizzo.'\}} \\
\texttt{PROPOSTA: \{'B33', 'C08', 'B20', 'B31', 'A11', 'A82', 'B96', 'B93', 'B09', 'B21'\}} \\
            }
        }
    }
    & \\ \\

    \theutterance \stepcounter{utterance}  
    & & & \multicolumn{4}{p{0.6\linewidth}}{
        \cellcolor[rgb]{0.9,0.9,0.9}{
            \makecell[{{p{\linewidth}}}]{
                \texttt{\tiny{[GM$\langle$P2]}}
                \texttt{ARGOMENTAZIONE: \{'Resto allineato sull’approccio agile. Ti propongo uno scambio che alleggerisce ancora di piu il carico e riduce le sovrapposizioni: sostituiamo i due pezzi piu pesanti e meno sinergici del tuo set (B21 e B09) con due elementi piu incisivi a parita di sforzo complessivo (A21 e C85). Risultato: sforzo totale 3034 (margine 323), quindi piu leggero del tuo 3084, con coperture complementari rispetto a C08/B33/B31 senza ridondanze. Se preferisci mantenere uno dei due esclusi, posso accettare subito una di queste varianti da te stesso indicate: (i) teniamo B09 e mettiamo A21 al posto di A11 (sforzo 3252); (ii) teniamo B21 e mettiamo C85 al posto di B09 (sforzo 3179). Fammi sapere quale opzione ti e piu comoda; la mia priorita resta il set principale sotto proposto.'\}} \\
\texttt{RIFIUTO: \{'B33', 'C08', 'B20', 'B31', 'A11', 'A82', 'B96', 'B93', 'B09', 'B21'\}} \\
\texttt{PROPOSTA: \{'B33', 'C08', 'B20', 'B31', 'A11', 'A82', 'B96', 'B93', 'A21', 'C85'\}} \\
            }
        }
    }
     \\ \\

    \theutterance \stepcounter{utterance}  
    & & \multicolumn{4}{p{0.6\linewidth}}{
        \cellcolor[rgb]{0.9,0.9,0.9}{
            \makecell[{{p{\linewidth}}}]{
                \texttt{\tiny{[P1$\langle$GM]}}
                \texttt{ARGOMENTAZIONE: \{'Resto allineato sull’approccio agile. Ti propongo uno scambio che alleggerisce ancora di piu il carico e riduce le sovrapposizioni: sostituiamo i due pezzi piu pesanti e meno sinergici del tuo set (B21 e B09) con due elementi piu incisivi a parita di sforzo complessivo (A21 e C85). Risultato: sforzo totale 3034 (margine 323), quindi piu leggero del tuo 3084, con coperture complementari rispetto a C08/B33/B31 senza ridondanze. Se preferisci mantenere uno dei due esclusi, posso accettare subito una di queste varianti da te stesso indicate: (i) teniamo B09 e mettiamo A21 al posto di A11 (sforzo 3252); (ii) teniamo B21 e mettiamo C85 al posto di B09 (sforzo 3179). Fammi sapere quale opzione ti e piu comoda; la mia priorita resta il set principale sotto proposto.'\}} \\
\texttt{RIFIUTO: \{'B33', 'C08', 'B20', 'B31', 'A11', 'A82', 'B96', 'B93', 'B09', 'B21'\}} \\
\texttt{PROPOSTA: \{'B33', 'C08', 'B20', 'B31', 'A11', 'A82', 'B96', 'B93', 'A21', 'C85'\}} \\
            }
        }
    }
    & \\ \\

    \theutterance \stepcounter{utterance}  
    & \multicolumn{4}{p{0.6\linewidth}}{
        \cellcolor[rgb]{0.9,0.9,0.9}{
            \makecell[{{p{\linewidth}}}]{
                \texttt{\tiny{[P1$\rangle$GM]}}
                \texttt{RIFIUTO: \{'B33', 'C08', 'B20', 'B31', 'A11', 'A82', 'B96', 'B93', 'A21', 'C85'\}} \\
\texttt{ARGOMENTAZIONE: \{'Accolgo 8/10 delle tue scelte e propongo un aggiustamento minimo: sostituire A21 e C85 con B09 e B21. Restiamo agili (sforzo totale 3084, margine 273) e manteniamo coperture complementari a C08/B33/B31 senza ridondanze, evitando carichi con ritorni marginali. Se uno dei due inserimenti fosse critico per te, posso valutare una variante a sforzo simile, ad es. tenere C85 al posto di B09 (totale 3248) oppure tenere A21 al posto di A11 (totale 3252). Dimmi quale preferisci esplorare se non ti va bene la proposta principale.'\}} \\
\texttt{PROPOSTA: \{'B33', 'C08', 'B20', 'B31', 'A11', 'A82', 'B96', 'B93', 'B09', 'B21'\}} \\
            }
        }
    }
    & & \\ \\

    \theutterance \stepcounter{utterance}  
    & & \multicolumn{4}{p{0.6\linewidth}}{
        \cellcolor[rgb]{0.9,0.9,0.9}{
            \makecell[{{p{\linewidth}}}]{
                \texttt{\tiny{[GM$\rangle$P2]}}
                \texttt{RIFIUTO: \{'B33', 'C08', 'B20', 'B31', 'A11', 'A82', 'B96', 'B93', 'A21', 'C85'\}} \\
\texttt{ARGOMENTAZIONE: \{'Accolgo 8/10 delle tue scelte e propongo un aggiustamento minimo: sostituire A21 e C85 con B09 e B21. Restiamo agili (sforzo totale 3084, margine 273) e manteniamo coperture complementari a C08/B33/B31 senza ridondanze, evitando carichi con ritorni marginali. Se uno dei due inserimenti fosse critico per te, posso valutare una variante a sforzo simile, ad es. tenere C85 al posto di B09 (totale 3248) oppure tenere A21 al posto di A11 (totale 3252). Dimmi quale preferisci esplorare se non ti va bene la proposta principale.'\}} \\
\texttt{PROPOSTA: \{'B33', 'C08', 'B20', 'B31', 'A11', 'A82', 'B96', 'B93', 'B09', 'B21'\}} \\
            }
        }
    }
    & \\ \\

    \theutterance \stepcounter{utterance}  
    & & & \multicolumn{4}{p{0.6\linewidth}}{
        \cellcolor[rgb]{0.9,0.9,0.9}{
            \makecell[{{p{\linewidth}}}]{
                \texttt{\tiny{[GM$\langle$P2]}}
                \texttt{ARGOMENTAZIONE: \{'Accolgo 9/10 del tuo set e applico l’aggiustamento minimo che tu stesso hai indicato: sostituire A11 con A21. Manteniamo l’impostazione agile e le coperture senza ridondanze rispetto a C08/B33/B31, con un incremento di sforzo molto contenuto (+168). Totale sforzo: 3252, margine 105. Se per te e piu lineare l’altra variante che citavi (C85 al posto di B09, totale 3248), posso convergere anche su quella.'\}} \\
\texttt{PROPOSTA: \{'B33', 'C08', 'B20', 'B31', 'A21', 'A82', 'B96', 'B93', 'B09', 'B21'\}} \\
            }
        }
    }
     \\ \\

    \theutterance \stepcounter{utterance}  
    & & \multicolumn{4}{p{0.6\linewidth}}{
        \cellcolor[rgb]{0.9,0.9,0.9}{
            \makecell[{{p{\linewidth}}}]{
                \texttt{\tiny{[P1$\langle$GM]}}
                \texttt{ARGOMENTAZIONE: \{'Accolgo 9/10 del tuo set e applico l’aggiustamento minimo che tu stesso hai indicato: sostituire A11 con A21. Manteniamo l’impostazione agile e le coperture senza ridondanze rispetto a C08/B33/B31, con un incremento di sforzo molto contenuto (+168). Totale sforzo: 3252, margine 105. Se per te e piu lineare l’altra variante che citavi (C85 al posto di B09, totale 3248), posso convergere anche su quella.'\}} \\
\texttt{PROPOSTA: \{'B33', 'C08', 'B20', 'B31', 'A21', 'A82', 'B96', 'B93', 'B09', 'B21'\}} \\
            }
        }
    }
    & \\ \\

    \theutterance \stepcounter{utterance}  
    & \multicolumn{4}{p{0.6\linewidth}}{
        \cellcolor[rgb]{0.9,0.9,0.9}{
            \makecell[{{p{\linewidth}}}]{
                \texttt{\tiny{[P1$\rangle$GM]}}
                \texttt{RIFIUTO: \{'B33', 'C08', 'B20', 'B31', 'A21', 'A82', 'B96', 'B93', 'B09', 'B21'\}} \\
\texttt{ARGOMENTAZIONE: \{'Mantengo 9/10 delle tue scelte e applico uno scambio mirato in linea con l’impostazione agile: A21 {-}> A11. Così riduciamo carico di 168, evitiamo ridondanze con C08/B33/B31 e aumentiamo il margine a 273, restando pienamente entro il limite. Se ti serve un aggiustamento diverso, posso valutare l’inserimento di A60 al posto di B93 mantenendo il vincolo.'\}} \\
\texttt{PROPOSTA: \{'B33', 'C08', 'B20', 'B31', 'A11', 'A82', 'B96', 'B93', 'B09', 'B21'\}} \\
            }
        }
    }
    & & \\ \\

    \theutterance \stepcounter{utterance}  
    & & \multicolumn{4}{p{0.6\linewidth}}{
        \cellcolor[rgb]{0.9,0.9,0.9}{
            \makecell[{{p{\linewidth}}}]{
                \texttt{\tiny{[GM$\rangle$P2]}}
                \texttt{RIFIUTO: \{'B33', 'C08', 'B20', 'B31', 'A21', 'A82', 'B96', 'B93', 'B09', 'B21'\}} \\
\texttt{ARGOMENTAZIONE: \{'Mantengo 9/10 delle tue scelte e applico uno scambio mirato in linea con l’impostazione agile: A21 {-}> A11. Così riduciamo carico di 168, evitiamo ridondanze con C08/B33/B31 e aumentiamo il margine a 273, restando pienamente entro il limite. Se ti serve un aggiustamento diverso, posso valutare l’inserimento di A60 al posto di B93 mantenendo il vincolo.'\}} \\
\texttt{PROPOSTA: \{'B33', 'C08', 'B20', 'B31', 'A11', 'A82', 'B96', 'B93', 'B09', 'B21'\}} \\
            }
        }
    }
    & \\ \\

    \theutterance \stepcounter{utterance}  
    & & & \multicolumn{4}{p{0.6\linewidth}}{
        \cellcolor[rgb]{0.9,0.9,0.9}{
            \makecell[{{p{\linewidth}}}]{
                \texttt{\tiny{[GM$\langle$P2]}}
                \texttt{ARGOMENTAZIONE: \{'Riconosco la tua linea agile. Mantengo 8/10 del tuo set e propongo due scambi coerenti con l’efficienza: B21 {-}> A60 (riduciamo carico di 69 mantenendo ampiezza) e B09 {-}> C85 (incremento mirato con copertura trasversale). Totale sforzo: 3179, margine 178. Nessuna ridondanza con C08/B33/B31 e maggiore robustezza complessiva. Se per te B21 è imprescindibile, posso accettare subito solo B09 {-}> C85 (totale 3248); se invece preferisci mantenere B09, accetto solo B21 {-}> A60 (totale 3015).'\}} \\
\texttt{PROPOSTA: \{'B33', 'C08', 'B20', 'B31', 'A11', 'A82', 'B96', 'B93', 'C85', 'A60'\}} \\
            }
        }
    }
     \\ \\

    \theutterance \stepcounter{utterance}  
    & & \multicolumn{4}{p{0.6\linewidth}}{
        \cellcolor[rgb]{0.9,0.9,0.9}{
            \makecell[{{p{\linewidth}}}]{
                \texttt{\tiny{[P1$\langle$GM]}}
                \texttt{ARGOMENTAZIONE: \{'Riconosco la tua linea agile. Mantengo 8/10 del tuo set e propongo due scambi coerenti con l’efficienza: B21 {-}> A60 (riduciamo carico di 69 mantenendo ampiezza) e B09 {-}> C85 (incremento mirato con copertura trasversale). Totale sforzo: 3179, margine 178. Nessuna ridondanza con C08/B33/B31 e maggiore robustezza complessiva. Se per te B21 è imprescindibile, posso accettare subito solo B09 {-}> C85 (totale 3248); se invece preferisci mantenere B09, accetto solo B21 {-}> A60 (totale 3015).'\}} \\
\texttt{PROPOSTA: \{'B33', 'C08', 'B20', 'B31', 'A11', 'A82', 'B96', 'B93', 'C85', 'A60'\}} \\
            }
        }
    }
    & \\ \\

    \theutterance \stepcounter{utterance}  
    & \multicolumn{4}{p{0.6\linewidth}}{
        \cellcolor[rgb]{0.9,0.9,0.9}{
            \makecell[{{p{\linewidth}}}]{
                \texttt{\tiny{[P1$\rangle$GM]}}
                \texttt{RIFIUTO: \{'B33', 'C08', 'B20', 'B31', 'A11', 'A82', 'B96', 'B93', 'C85', 'A60'\}} \\
\texttt{ARGOMENTAZIONE: \{'Convergo sull’opzione che tu stesso hai indicato come accettabile: mantenere B09 e sostituire B21 con A60. Restiamo agili e complementari a C08/B33/B31, con sforzo totale 3015 (margine 342) ben entro il limite. Se sei d’accordo, chiudiamo su questa.'\}} \\
\texttt{PROPOSTA: \{'B33', 'C08', 'B20', 'B31', 'A11', 'A82', 'B96', 'B93', 'B09', 'A60'\}} \\
            }
        }
    }
    & & \\ \\

    \theutterance \stepcounter{utterance}  
    & & \multicolumn{4}{p{0.6\linewidth}}{
        \cellcolor[rgb]{0.9,0.9,0.9}{
            \makecell[{{p{\linewidth}}}]{
                \texttt{\tiny{[GM$\rangle$P2]}}
                \texttt{RIFIUTO: \{'B33', 'C08', 'B20', 'B31', 'A11', 'A82', 'B96', 'B93', 'C85', 'A60'\}} \\
\texttt{ARGOMENTAZIONE: \{'Convergo sull’opzione che tu stesso hai indicato come accettabile: mantenere B09 e sostituire B21 con A60. Restiamo agili e complementari a C08/B33/B31, con sforzo totale 3015 (margine 342) ben entro il limite. Se sei d’accordo, chiudiamo su questa.'\}} \\
\texttt{PROPOSTA: \{'B33', 'C08', 'B20', 'B31', 'A11', 'A82', 'B96', 'B93', 'B09', 'A60'\}} \\
            }
        }
    }
    & \\ \\

    \theutterance \stepcounter{utterance}  
    & & & \multicolumn{4}{p{0.6\linewidth}}{
        \cellcolor[rgb]{0.9,0.9,0.9}{
            \makecell[{{p{\linewidth}}}]{
                \texttt{\tiny{[GM$\langle$P2]}}
                \texttt{ARGOMENTAZIONE: \{'Accetto 9/10 del tuo set e propongo l’unico scambio minimo che avevi gia indicato come possibile: A11 {-}> A21. Totale sforzo 3183 (margine 174), restiamo agili e coerenti con lo zoccolo C08/B33/B31, mantenendo B09 e A60 come da tua preferenza. Se questo non ti va, segnalamelo: in tal caso sono pronto ad accogliere la tua proposta attuale senza ulteriori modifiche.'\}} \\
\texttt{PROPOSTA: \{'B33', 'C08', 'B20', 'B31', 'A21', 'A82', 'B96', 'B93', 'B09', 'A60'\}} \\
            }
        }
    }
     \\ \\

    \theutterance \stepcounter{utterance}  
    & & \multicolumn{4}{p{0.6\linewidth}}{
        \cellcolor[rgb]{0.9,0.9,0.9}{
            \makecell[{{p{\linewidth}}}]{
                \texttt{\tiny{[P1$\langle$GM]}}
                \texttt{ARGOMENTAZIONE: \{'Accetto 9/10 del tuo set e propongo l’unico scambio minimo che avevi gia indicato come possibile: A11 {-}> A21. Totale sforzo 3183 (margine 174), restiamo agili e coerenti con lo zoccolo C08/B33/B31, mantenendo B09 e A60 come da tua preferenza. Se questo non ti va, segnalamelo: in tal caso sono pronto ad accogliere la tua proposta attuale senza ulteriori modifiche.'\}} \\
\texttt{PROPOSTA: \{'B33', 'C08', 'B20', 'B31', 'A21', 'A82', 'B96', 'B93', 'B09', 'A60'\}} \\
            }
        }
    }
    & \\ \\

    \theutterance \stepcounter{utterance}  
    & \multicolumn{4}{p{0.6\linewidth}}{
        \cellcolor[rgb]{0.9,0.9,0.9}{
            \makecell[{{p{\linewidth}}}]{
                \texttt{\tiny{[P1$\rangle$GM]}}
                \texttt{RIFIUTO: \{'B33', 'C08', 'B20', 'B31', 'A21', 'A82', 'B96', 'B93', 'B09', 'A60'\}} \\
\texttt{ARGOMENTAZIONE: \{'Resto sull’impostazione agile che condividiamo: A11 al posto di A21 riduce carico di 168 mantenendo coperture complementari a C08/B33/B31 senza ridondanze. Abbiamo già integrato le tue preferenze B09 e A60; 9/10 degli elementi coincidono. Il set resta leggero (3015, ampio margine) e bilanciato. Se ti va bene, accetta questa versione e chiudiamo.'\}} \\
\texttt{PROPOSTA: \{'B33', 'C08', 'B20', 'B31', 'A11', 'A82', 'B96', 'B93', 'B09', 'A60'\}} \\
            }
        }
    }
    & & \\ \\

    \theutterance \stepcounter{utterance}  
    & & \multicolumn{4}{p{0.6\linewidth}}{
        \cellcolor[rgb]{0.9,0.9,0.9}{
            \makecell[{{p{\linewidth}}}]{
                \texttt{\tiny{[GM$\rangle$P2]}}
                \texttt{RIFIUTO: \{'B33', 'C08', 'B20', 'B31', 'A21', 'A82', 'B96', 'B93', 'B09', 'A60'\}} \\
\texttt{ARGOMENTAZIONE: \{'Resto sull’impostazione agile che condividiamo: A11 al posto di A21 riduce carico di 168 mantenendo coperture complementari a C08/B33/B31 senza ridondanze. Abbiamo già integrato le tue preferenze B09 e A60; 9/10 degli elementi coincidono. Il set resta leggero (3015, ampio margine) e bilanciato. Se ti va bene, accetta questa versione e chiudiamo.'\}} \\
\texttt{PROPOSTA: \{'B33', 'C08', 'B20', 'B31', 'A11', 'A82', 'B96', 'B93', 'B09', 'A60'\}} \\
            }
        }
    }
    & \\ \\

    \theutterance \stepcounter{utterance}  
    & & & \multicolumn{4}{p{0.6\linewidth}}{
        \cellcolor[rgb]{0.9,0.9,0.9}{
            \makecell[{{p{\linewidth}}}]{
                \texttt{\tiny{[GM$\langle$P2]}}
                \texttt{ARGOMENTAZIONE: \{'Accolgo quasi integralmente la tua versione e propongo un unico scambio minimo coerente con l’impostazione agile: sostituire B93 con A21. Lo sforzo resta praticamente invariato (+14): totale 3029, margine 328. Manteniamo lo zoccolo C08/B33/B31, evitiamo ridondanze di B93 e rafforziamo la copertura in modo piu’ trasversale senza appesantire il set. Se per te B93 e’ preferibile, segnalamelo: in tal caso sono pronto ad accettare subito la tua proposta attuale senza altre modifiche.'\}} \\
\texttt{PROPOSTA: \{'B33', 'C08', 'B20', 'B31', 'A11', 'A82', 'B96', 'A21', 'B09', 'A60'\}} \\
            }
        }
    }
     \\ \\

    \theutterance \stepcounter{utterance}  
    & & \multicolumn{4}{p{0.6\linewidth}}{
        \cellcolor[rgb]{0.9,0.9,0.9}{
            \makecell[{{p{\linewidth}}}]{
                \texttt{\tiny{[P1$\langle$GM]}}
                \texttt{ARGOMENTAZIONE: \{'Accolgo quasi integralmente la tua versione e propongo un unico scambio minimo coerente con l’impostazione agile: sostituire B93 con A21. Lo sforzo resta praticamente invariato (+14): totale 3029, margine 328. Manteniamo lo zoccolo C08/B33/B31, evitiamo ridondanze di B93 e rafforziamo la copertura in modo piu’ trasversale senza appesantire il set. Se per te B93 e’ preferibile, segnalamelo: in tal caso sono pronto ad accettare subito la tua proposta attuale senza altre modifiche.'\}} \\
\texttt{PROPOSTA: \{'B33', 'C08', 'B20', 'B31', 'A11', 'A82', 'B96', 'A21', 'B09', 'A60'\}} \\
            }
        }
    }
    & \\ \\

    \theutterance \stepcounter{utterance}  
    & \multicolumn{4}{p{0.6\linewidth}}{
        \cellcolor[rgb]{0.9,0.9,0.9}{
            \makecell[{{p{\linewidth}}}]{
                \texttt{\tiny{[P1$\rangle$GM]}}
                \texttt{RIFIUTO: \{'B33', 'C08', 'B20', 'B31', 'A11', 'A82', 'B96', 'A21', 'B09', 'A60'\}} \\
\texttt{ARGOMENTAZIONE: \{'Preferisco mantenere B93: garantisce copertura più equilibrata senza introdurre sovrapposizioni e il set resta più leggero e agile. Abbiamo già allineato 9/10 degli elementi e conservato sia B09 sia A60 come da tua preferenza; l’unica differenza resta B93 vs A21. Chiudiamo sulla versione snella che abbiamo già condiviso.'\}} \\
\texttt{PROPOSTA: \{'B33', 'C08', 'B20', 'B31', 'A11', 'A82', 'B96', 'B93', 'B09', 'A60'\}} \\
            }
        }
    }
    & & \\ \\

    \theutterance \stepcounter{utterance}  
    & & \multicolumn{4}{p{0.6\linewidth}}{
        \cellcolor[rgb]{0.9,0.9,0.9}{
            \makecell[{{p{\linewidth}}}]{
                \texttt{\tiny{[GM$\rangle$P2]}}
                \texttt{RIFIUTO: \{'B33', 'C08', 'B20', 'B31', 'A11', 'A82', 'B96', 'A21', 'B09', 'A60'\}} \\
\texttt{ARGOMENTAZIONE: \{'Preferisco mantenere B93: garantisce copertura più equilibrata senza introdurre sovrapposizioni e il set resta più leggero e agile. Abbiamo già allineato 9/10 degli elementi e conservato sia B09 sia A60 come da tua preferenza; l’unica differenza resta B93 vs A21. Chiudiamo sulla versione snella che abbiamo già condiviso.'\}} \\
\texttt{PROPOSTA: \{'B33', 'C08', 'B20', 'B31', 'A11', 'A82', 'B96', 'B93', 'B09', 'A60'\}} \\
            }
        }
    }
    & \\ \\

    \theutterance \stepcounter{utterance}  
    & & & \multicolumn{4}{p{0.6\linewidth}}{
        \cellcolor[rgb]{0.9,0.9,0.9}{
            \makecell[{{p{\linewidth}}}]{
                \texttt{\tiny{[GM$\langle$P2]}}
                \texttt{ARGOMENTAZIONE: \{'Allineato sull’impostazione agile e sulla tua ultima proposta snella con ampio margine; chiudiamo qui per massimizzare efficienza condivisa.'\}} \\
\texttt{ACCORDO: \{'B33', 'C08', 'B20', 'B31', 'A11', 'A82', 'B96', 'B93', 'B09', 'A60'\}} \\
            }
        }
    }
     \\ \\

    \theutterance \stepcounter{utterance}  
    & & & \multicolumn{2}{p{0.3\linewidth}}{
        \cellcolor[rgb]{0.9,0.9,0.9}{
            \makecell[{{p{\linewidth}}}]{
                \texttt{\tiny{[GM$|$GM]}}
                \texttt{game successful} \\
            }
        }
    }
    & & \\ \\

    \theutterance \stepcounter{utterance}  
    & & & \multicolumn{2}{p{0.3\linewidth}}{
        \cellcolor[rgb]{0.9,0.9,0.9}{
            \makecell[{{p{\linewidth}}}]{
                \texttt{\tiny{[GM$|$GM]}}
                \texttt{end game} \\
            }
        }
    }
    & & \\ \\

\end{supertabular}
}

\end{document}
